\documentclass[a4paper, 11pt]{article}
\usepackage[left=3cm, top=3cm, right=2cm, bottom=2cm]{geometry}
%pacote que define as margens
\usepackage{helvet} % carrega o documento com Fonte ARIAL
\renewcommand{\familydefault}{\sfdefault} % define a fonte ARIAL como padrão do texto
\usepackage[utf8]{inputenc} %pacote que normaliza a acentuação
\usepackage[T1]{fontenc}
\usepackage[alf,abnt-repeated-title-omit=yes,abnt-emphasize=bf,abnt-etal-list=0]{abntex2cite} %define os parâmetros necessários para as referências e citações de acordo com a ABNT
\usepackage[brazil]{babel} %pacote que define o idioma
\usepackage{graphics} %pacote que permite a inserção de imagens
\usepackage{graphicx} %pacote que permite a inserção de imagens
\pagestyle{myheadings} %pacote que define o cabeçalho com o número de página no canto superior direito
\usepackage{setspace} %pacote que permite especificar o espaçamento entre linhas no documento
\usepackage{indentfirst} %pacote que faz indentação na primeira linha do parágrafo
\setlength{\parindent}{1.25cm} %pacote que define o tamanho da indentação
\usepackage{float} %Pacote que permite tabelas flutuarem em qualquer posição
\usepackage{amsthm} % Pacote que permite a utilização de teoremas
\usepackage{amsmath} % Pacote que permite a utilização de fontes matemáticas
\usepackage{amsfonts} % Pacote que permite a utilização de fontes matemáticas
\usepackage{subcaption} % Pacote que permite que imagens possam ser inseridas em composições de imagens lado a lado com o \subfigure
\usepackage{amssymb} %4 Pacotes responsáveis pela formatação matemática
\usepackage{rotating} %permite que textos, figuras e tabelas possam ser rotacionados
\usepackage{array}    % permite e vetorização de tabelas
\usepackage{longtable} % permite a elaboração de tabelas complexas
\usepackage{tikz} % criação de objetos vetoriais
\usetikzlibrary{shapes, arrows.meta, positioning}
\usepackage{pdflscape} % permite alterar o estilo da pagina entre portrato e paisagem
\usepackage{hyperref}
\usepackage{sectsty}
\usepackage[most]{tcolorbox}
\usepackage{listings}        % Para listar códigos
\usepackage[acronym]{glossaries} % cria glossário de siglas
\usepackage{url}
\usepackage[table,xcdraw]{xcolor}
\usepackage{colortbl}
\usepackage{cite}
\hypersetup{
    colorlinks=true,
    linkcolor=darkblue,
    filecolor=magenta,
    urlcolor=blue,
    citecolor=black,
    pdftitle={mybib},
    pdfpagemode=FullScreen
}

\sectionfont{\fontsize{12}{12}\selectfont}
\subsectionfont{\fontsize{12}{12}\selectfont}
\subsubsectionfont{\fontsize{12}{12}\selectfont}

%\usepackage[brazilian,hyperpageref]{backref}	% Paginas com as citações na bibl



%\usepackage{pgfgantt}  % gera gantt em latex
%\makeglossaries


\lstset{
    inputencoding=utf8,
    extendedchars=true,
    literate={á}{{\'a}}1 {é}{{\'e}}1 {í}{{\'i}}1 {ó}{{\'o}}1 {ú}{{\'u}}1
        {Á}{{\'A}}1 {É}{{\'E}}1 {Í}{{\'I}}1 {Ó}{{\'O}}1 {Ú}{{\'U}}1
        {à}{{\`a}}1 {è}{{\`e}}1 {ì}{{\`i}}1 {ò}{{\`o}}1 {ù}{{\`u}}1
        {À}{{\`A}}1 {È}{{\`E}}1 {Ì}{{\`I}}1 {Ò}{{\`O}}1 {Ù}{{\`U}}1
        {ã}{{\~a}}1 {õ}{{\~o}}1 {Ã}{{\~A}}1 {Õ}{{\~O}}1
        {â}{{\^a}}1 {ê}{{\^e}}1 {î}{{\^i}}1 {ô}{{\^o}}1 {û}{{\^u}}1
        {Â}{{\^A}}1 {Ê}{{\^E}}1 {Î}{{\^I}}1 {Ô}{{\^O}}1 {Û}{{\^U}}1
        {ç}{{\c{c}}}1 {Ç}{{\c{C}}}1
}

\definecolor{codegreen}{rgb}{0,0.6,0}
\definecolor{codegray}{rgb}{0.5,0.5,0.5}
\definecolor{codepurple}{rgb}{0.58,0,0.82}
\definecolor{backcolour}{rgb}{0.95,0.95,0.92}

\lstdefinestyle{mystyle}{
    backgroundcolor=\color{backcolour},
    commentstyle=\color{codegreen},
    keywordstyle=\color{magenta},
    numberstyle=\tiny\color{codegray},
    stringstyle=\color{codepurple},
    basicstyle=\ttfamily\footnotesize,
    breakatwhitespace=false,
    breaklines=true,
    captionpos=b,
    keepspaces=true,
    numbers=left,
    numbersep=5pt,
    showspaces=false,
    showstringspaces=false,
    showtabs=false,
    tabsize=2
}

\lstset{style=mystyle}

% Document
\begin{document}
    
\section*{Lista de Exercício – Seleção de Materiais - Prova P1}

\begin{itemize}
    \item \textbf{Professor}: Dr. Márcio Wagner
    \item \textbf{Disciplina}: Seleção de Materiais
    \item \textbf{Data de Entrega}: 12/02/2025
    \item \textbf{Aluno}: Alan Henrique Pereira Miranda
    \item \textbf{Matrícula}: 202102140072
\end{itemize}

\section{Questão 1}
\subsection*{a) Fluxograma de Projeto}
O fluxograma de projeto é uma representação visual das etapas envolvidas no desenvolvimento de um produto ou sistema. Ele inclui fases como identificação do problema, definição de requisitos, concepção, seleção de materiais, modelagem, prototipagem e testes.

\subsection*{b) Sistema Técnico e suas Subdivisões}
Um sistema técnico é um conjunto de elementos que trabalham juntos para realizar uma função específica. Suas subdivisões incluem:
\begin{itemize}
    \item \textbf{Entrada}: matérias-primas, informação, energia.
    \item \textbf{Processamento}: transformação de insumos.
    \item \textbf{Saída}: produto final, serviço ou energia gerada.
    \item \textbf{Controle}: ajustes e feedbacks para manutenção do sistema.
\end{itemize}

\subsection*{c) Estrutura da Função}
Define a relação entre os elementos de um sistema e como esses elementos interagem para cumprir a função desejada. É composta por funções principais e secundárias.

\subsection*{d) Caminho Convoluto do Projeto}
Refere-se às iterativas e complexas interações entre diferentes etapas do projeto, que podem exigir revisões contínuas.

\subsection{*e) Problema Central de Seleção de Materiais}
A seleção envolve a escolha do material ideal considerando função, processo de fabricação, forma e custo.

\section{Famílias de Materiais e seus Principais Membros}
\begin{itemize}
    \item \textbf{Metais}: Aço, alumínio, cobre, titânio.
    \item \textbf{Cerâmicas}: Alumina, zircônia, carbeto de silício.
    \item \textbf{Polímeros}: Polietileno, polipropileno, nylon.
    \item \textbf{Compósitos}: Fibra de carbono, fibra de vidro.
    \item \textbf{Elastômeros}: Borracha natural, neoprene, silicone.
\end{itemize}

\section{Propriedades dos Materiais}
\begin{itemize}
    \item \textbf{Cerâmicas}: Alta dureza, baixa tenacidade, resistência ao calor.
    \item \textbf{Vidros}: Transparência, fragilidade, resistência a ataques químicos.
    \item \textbf{Polímeros}: Flexibilidade, baixa densidade, boa moldabilidade.
    \item \textbf{Elastômeros}: Alta elasticidade, resistência ao impacto.
    \item \textbf{Metais}: Alta resistência mecânica, ductilidade, condutividade térmica.
    \item \textbf{Híbridos}: Combina propriedades de diferentes materiais, como rigidez e leveza.
\end{itemize}

\section{Tipos de Informações Necessárias aos Projetos}
\begin{itemize}
    \item Propriedades mecânicas e físicas.
    \item Processabilidade.
    \item Durabilidade e custos.
\end{itemize}

\section{Materiais Quanto à Ductilidade}
\subsection*{a) Dúcteis}
Sofrem grandes deformações plásticas antes da fratura.

\subsection*{b) Polímeros Dúcteis}
Apresentam deformabilidade acentuada, podendo esticar significativamente antes de romper.

\subsection*{c) Cerâmicas Frágeis}
Rompem com pouca deformação plástica.

\section{O que é Limite de Fadiga?}
A tensão abaixo da qual um material pode suportar um número infinito de ciclos de carga sem falhar.

\section{O que é Dureza?}
Resistência à penetração ou ao desgaste. Principais escalas: Brinell, Rockwell, Vickers.

\section{O que é Tenacidade à Fratura?}
Capacidade de um material resistir à propagação de trincas.

\section{Condutividade Térmica}
Medida da capacidade de transferência de calor. Governada pela equação de Fourier.

\section{Expansão Térmica}
Variação dimensional de um material com a temperatura, descrita pelas equações:

\begin{itemize}
    \item \textbf{Linear}: \( \Delta L = \alpha L_0 \Delta T \)
    \item \textbf{Superficial}: \( \Delta A = 2\alpha A_0 \Delta T \)
    \item \textbf{Volumétrica}: \( \Delta V = 3\alpha V_0 \Delta T \)
\end{itemize}

\section{Como são obtidos os valores de resistência dos materiais?}
Por meio de ensaios mecânicos como tração, compressão e impacto.

\section{Diagrama de Barras Únicas}
Representação gráfica de uma propriedade de vários materiais.

\section{Diagramas de Propriedades Múltiplas}
Gráfico com duas propriedades distintas para seleção de materiais.

\section{Explicação dos Diagramas}
\begin{itemize}
    \item \( E-\rho \) (Módulo de Young vs. Densidade)
    \item \( R-\rho \) (Resistência vs. Densidade)
    \item \( E^*/R^* \) (Módulo Específico vs. Resistência Específica)
    \item Tenacidade vs. Módulo
    \item Tenacidade à Fratura vs. Resistência
\end{itemize}

\section{Por que Conectar Material e Função?}
Para garantir que o material atenda às demandas do projeto com eficiência.

\section{Estratégia para Seleção de Materiais}
Definir requisitos, classificar opções e otimizar a escolha.

\section{Passos: Tradução, Triagem, Classificação e Documentação}
Transformar requisitos do projeto em critérios técnicos e realizar a seleção.

\section{O que é Função, Restrição, Objetivo e Variável Livre?}
\begin{itemize}
    \item \textbf{Função}: Propósito do componente.
    \item \textbf{Restrição}: Limites obrigatórios.
    \item \textbf{Objetivo}: O que otimizar.
    \item \textbf{Variável Livre}: Parâmetros ajustáveis.
\end{itemize}

\end{document}